\chapterstart{第一阶段：前沿突破}
\label{sec:frontier}

第一阶段的任务是在严格的数据预算下训练出具有竞争力的语言模型。Qiushi Engine 需要同时解决三个问题：有限文本应怎样安排，模型应怎样构造，以及训练接近预算上限时还能怎样改进。有效方案在文本处理、结构设计和完整评测的反复比较中形成。它最终留下的不只是一组高分权重，也是一套可以继续训练、比较和分析的实验基础。

\subsection{联合选择表示、结构与学习方式}

模型设计首先要决定怎样把文本变成可学习的对象。字符表示把文字拆得更细，子词表示把常用片段合并；前者需要处理更长序列，后者需要为较多符号学习表示。预测方式也有区别：自回归模型根据前文预测后文，掩码模型则遮盖部分词，再依据未被遮盖的上下文恢复它们。本研究比较了这些选择，以及模型深度、宽度、词表和训练安排的组合。

其他探索包括只激活部分计算分支的稀疏结构，以及给模型增加额外记忆或词面特征。它们有的降低了某项训练损失，但最终需要回答同一个问题：在可比较的预算下，是否提高了所需语言能力？因此，结构比较需要对齐共同参数的初值，分词比较需要重新统计序列和预算，学习目标比较需要同时观察受益与受损任务。研究逐渐从单个部件的局部效果，转向能够共同工作的完整训练方案。

最终前沿路线收束于三项能够共同使用的设计：以紧凑重述重新分配文本预算，以残差结构保留可延续的学习空间，再以完整评测选择晚期增量。其他结构试验仍保留各自结果，详见第\ref{sec:lineage}章及附录\ref{app:catalog}。以下沿最终模型的实际形成过程说明这三项设计。

\subsection{紧凑重述与预算再分配}

训练数据包含原文及相应改写。本文把原文称为“来源”，把保留对应内容的另一种表达称为“重述”；两者组成一个来源—重述对。相关实验材料也将重述称为文本视图（text view），指同一内容的另一种文字表达。第二种表达提供了不同于逐字重复的预测任务，同时也占用词数预算。紧凑重述在保留内容联系的同时减少冗余，将节约的词数用于更多可学习材料。

具体做法是缩短配对文本中的重述部分，再将释放的预算投入更多段落对。核心段落对由 10,094 增至 12,155，增加 2,061 对，约为 20.42\%。这项改动发生在语料池内一个大小固定的文本子集，以下称为替换数据块。该块总计 423,520 词，由 261,803 词来源文本、161,708 词紧凑重述及 9 词补齐组成，涉及 4,529 个来源文档；同一文档可以贡献多个段落对。

\begin{table}[htbp]
\caption{固定词数下的紧凑重述与经验扩展。压缩重述部分后，替换数据块容纳了更多来源—重述对，总词数保持为 423,520。}
\label{tab:compact}
\begin{tabularx}{\linewidth}{@{}Yrr@{}}
\toprule 比较对象 & 数量 & 研究含义\\\midrule
原有核心段落对 & 10,094 & 配对学习基础\\
新增段落对 & 2,061 & 压缩后重新投入的经验\\
最终段落对 & 12,155 & 较原有数量增加 20.42\%\\
来源文档 & 4,529 & 一个文档可提供多个段落对\\
来源文本／紧凑重述（词） & 261,803／161,708 & 保留两种表达的联系\\
替换数据块总词数 & 423,520 & 包含 9 词补齐\\
\bottomrule
\end{tabularx}
\end{table}

这项设计同时改变了表达形式和经验覆盖：重述压缩调整了模型需要预测的内容，新增段落对则让同等预算容纳更多来源信息。分开辨明这两种作用，后来成为固定预算替换与目标删除实验的重要问题。

训练中还出现了值得解释的排序变化。为较快比较候选，研究采用 BLiMP、Supplement、EWoK、Entity Tracking、COMPS、GlobalPIQA 与 Reading 的七项均值。这一研究指标包含 Reading、不含 (Super)GLUE，与官方 NLP Average 的七项组成不同。在相应匹配比较中，紧凑重述方案相对参照的差值，在累计 20M 词时约为 $-0.83$，到 70M 与 80M 词时分别为 $+1.29$ 和 $+1.35$。早期劣势转为晚期优势，促使研究把学习阶段纳入数据价值的判断：一种经验可能要在模型形成一定基础能力后，才带来明显收益。

本阶段还研究了连续窗口装填，即把相邻文本按模型可处理的长度组织为训练输入。这类探索积累了词数利用、截断处理和上下文组织方法。最终模型继承了其中明确采用的数据和权重，其他路线则作为独立工程成果保留。

\subsection{残差模型与可延续的增量学习}

代表模型采用 DeBERTa-v2 风格的掩码语言建模架构\citep{he2021deberta,microsoft2021debertav2}：8 层、隐藏维度 480、8 个注意力头，词表大小为 16,384。隐藏维度是每个位置的内部向量长度；注意力头并行学习从哪些上下文位置提取信息。预测窗口是模型一次处理时可以共同访问的文本范围，主模型预训练采用 256 个 token。分词使用字节级 BPE（Byte Pair Encoding），通过合并常见字节片段构造词表，词表由受限语料训练得到。基本训练采用 15\% 的整词遮盖：一个词被选中时，其对应的子词一起遮盖。

残差结构让一条较小的计算支路对原表示作补充。可以把它理解为保留已有计算，再学习一项修正；这里先把 480 维表示压到 128 维，经变换后映回原空间。因此，研究能够控制新增计算的大小，并在后续训练中区分原有参数与新增参数。

在主路径联合训练阶段，残差支路可概括为
\begin{equation}
 h'=h+s\,U\,\sigma\!\left(D\,\operatorname{Norm}(h)\right),
\label{eq:residual}
\end{equation}
其中 $h$ 是原表示，$D$ 与 $U$ 分别压缩和恢复维度，$\sigma$ 是非线性变换，$\operatorname{Norm}$ 为归一化，$s$ 控制补充量。上投影采用零初始化，使新支路在起点不改变原表示；主干与支路随后共同训练。主模型联合训练采用 $s=1.75$，从头学习，不载入外部预训练权重。该构造借鉴了残差与参数高效学习方法\citep{bachlechner2021rezero,houlsby2019adapters}；本研究将其与有限预算下的数据组织、从头训练及后续增量学习结合。

主模型在累计 82,012,495 词曝光时具有 35,463,008 个参数。随后冻结已有参数，增加一条专用的可训练增量分支：共 995,584 个参数、48 个参数张量，最终总参数数为 36,458,592。该阶段使用 101 次参数更新，计入 3,992,800 词曝光，累计达到 86,005,295 词。

\begin{table}[htbp]
\caption{第一代代表模型的主要训练层次。两行属于实际模型谱系，冻结比例不直接代表相同比例的计算节约。}
\label{tab:lineage}
\small
\begin{tabularx}{\linewidth}{@{}P{32mm}Yrr@{}}
\toprule 阶段 & 参数更新方式 & 总参数数 & 累计词曝光\\\midrule
主模型联合训练 & 主模型与残差支路共同学习 & 35,463,008 & 82,012,495\\
冻结基座增量训练 & 仅更新新增分支 & 36,458,592 & 86,005,295\\
\bottomrule
\end{tabularx}
\end{table}

增量训练同时使用普通预测损失与 KL 保持项。普通损失训练模型正确预测文本；KL 散度（Kullback--Leibler divergence）衡量两个模型对候选词的预测概率分布有多大差异，保持项通过限制这一差异来约束原有预测的变化。第一代已经采用了这一能力保持设计。后续研究在此基础上继续追问：应在什么输入上保持哪些功能，怎样兼顾需要获得的新能力？

专用分支提供了明确的控制方式：关闭新增路径时恢复冻结基座的输出，开启时加入学到的修正。研究在开启状态下评价新旧能力能否共存，并单独记录可训练参数与计算成本。可训练参数约占总量的 2.73\%。冻结参数省去了相应参数的更新，但为了训练分布在多层中的增量模块，梯度仍需通过其上方的冻结运算传回；此外还需计算参照模型的输出。因此，参数更新量与实际计算量分别统计。

\subsection{连贯关系与晚期学习}

冻结基座、只训练增量分支，使研究能够在共同模型基础上比较文本组织。一个关键对照固定起点、增量结构、词曝光和更新数，只打散训练行内的连贯片段。残差尺度为 1.0 时，连贯方案的七项均值为 44.1064，打散方案为 43.1214，差值约 0.9850。相同词语因连贯组织不同而产生明确差异，由此将研究问题推进到文本关系本身。

\begin{table}[htbp]
\caption{晚期学习的参照与对照。七项均值由 BLiMP、Supplement、EWoK、Entity、COMPS、GlobalPIQA 与 Reading 构成。只有最后两行是这里最接近的组织方式对照；普通续训与增量训练还改变了更新参数、优化器和目标。}
\label{tab:late}
\small
\begin{tabularx}{\linewidth}{@{}Yrr@{}}
\toprule 条件 & 累计词曝光 & 七项均值\\\midrule
续训前主模型 & 82,012,495 & 43.9594\\
普通联合续训 & 86,006,729 & 43.7707\\
对应关系打乱的增量训练 & 86,005,413 & 44.0129\\
连贯增量训练（尺度 1.0） & 86,005,295 & 44.1064\\
行内片段打散的增量训练 & 86,005,295 & 43.1214\\
\bottomrule
\end{tabularx}
\end{table}

新增分支在尺度 1.0 下训练，随后比较不同的残差尺度，选择 0.75 作为发布配置。尺度直接乘在增量分支的输出上，0.75 表示向原表示加入该分支输出的四分之三。这是依据评测作出的模型选择。最终第一代的本地完整 Overall 为 42.0240，公开显示为 42.02。

综合指标与逐题行为的变化，引出了增量学习如何兼顾新旧能力的问题。

\subsection{第一阶段留下的科学起点}

第一阶段留下了可续训的模型、来源—重述材料与三个关键现象：连贯组织优于片段打散，紧凑重述的优势随学习阶段变化，局部改善未必带来整体收益。Qiushi Engine 因而进一步研究关系如何被学会、被新输入调用并在后续训练中保持；配对文本、共同基座和完整评测构成了机制实验的基础。
